Basis data didefinisikan sebagai kumpulan informasi yang terorganisir dan saling terkait. Tujuan utama basis data adalah untuk mengelola sejumlah besar data yang disimpan dan menyediakan informasi darinya~\cite{database2022jagdish}. Desain basis data adalah proses menghasilkan model data yang terperinci dari sebuah basis data~\cite{database2019date}. Model data \textit{entity-relationship} (E-R) adalah model data yang banyak digunakan untuk desain basis data. Model ini menyediakan representasi grafis yang nyaman untuk melihat data, hubungan, dan batasan~\cite{database2019abraham}. Menurut Conolly dan Begg~\cite{conolly2015database}, terdapat tiga fase dalam desain basis data yaitu \textit{conceptual database design, logical database design,} dan \textit{physical database design}.


\subsection{\textit{Conceptual Database Design}}
\textit{Conceptual data design} merupakan tahap untuk membangun representasi konseptual dari basis data, yang meliputi identifikasi entitas, hubungan, dan atribut penting~\cite{conolly2015database}. Perancangan \textit{conceptual data modeling} akan dilakukan dengan mengindentifikasi \textit{entity} berdasarkan hasil wawancara dengan \textit{Manager Human Resouces} di PT Password Solusi Sistem. Pada \textbf{\ref{tab:indentifikasi-entity}} menunjukkan hasil identifikasi \textit{entity} berdasarkan kebutuhan MyPassword. Pada \textbf{\ref{tab:indentifikasi-hubungan-antar-entity}} menunjukkan hubungan antar \textit{entity}. \textit{Conceptual data modeling} dalam bentuk diagram dapat dilihat pada \hyperref[sec:conceptual-database-design]{\textbf{Lampiran~VIII}}.
%%%%%%%%%%%%%%%%%%%%%%%%%%%%%%%%%%%%%%%%%%%%%%%%%%%%%%%%%%%%%%%%%
\FloatBarrier
\begin{longtable}{ |p{3cm}|p{6cm}|p{6cm}|}
\caption{Identifikasi \textit{Entity}}\label{tab:indentifikasi-entity} 
\\
\hline
\textbf{\textit{Entity}}
&\textbf{Deskripsi}
&\textbf{Kejadian / Kemunculan}
\\ 
\hline

\textit{User}
&Informasi tentang \textit{user} pada aplikasi presensi \textit{online}.
&Data \textit{user} akan dibuat ketika \textit{Staff} atau \textit{Manager Infrastructure} melakukan registrasi.
\\ 
\hline

\textit{Shift}
&Informasi tentang jenis \textit{shift} kerja.
&Data \textit{Shift} ditambahkan oleh \textit{Manager Human Resources} dan \textit{Manager Infrastructure}.
\\ 
\hline

\textit{ShiftDay}
&Informasi tentang hari dan waktu spesifik dari setiap jenis \textit{shift} yang akan digunakan untuk perhitungan lembur.
&Data \textit{ShiftDay} ditambahkan oleh \textit{Manager Human Resources} dan \textit{Manager Infrastructure}.
\\ 
\hline

\textit{Holiday}
&Informasi tentang jenis libur perusahaan.
&Data \textit{Holiday} ditambahkan oleh \textit{Manager Human Resources} dan \textit{Manager Infrastructure}.
\\ 
\hline

\textit{HolidayDay}
&Informasi tentang hari dan waktu spesifik dari setiap jenis \textit{holiday} yang akan digunakan untuk perhitungan lembur.
&Data \textit{HolidayDay} ditambahkan oleh \textit{Manager Human Resources} dan \textit{Manager Infrastructure}.
\\ 
\hline

\textit{Department}
&Informasi tentang divisi dari setiap \textit{user}.
&Data \textit{Department} ditambahkan oleh \textit{Manager Human Resources} dan \textit{Manager Infrastructure}.
\\ 
\hline

\textit{Role}
&Informasi tentang jabatan dari setiap \textit{user}.
&Data \textit{Role} ditambahkan oleh \textit{Manager Human Resources} dan \textit{Manager Infrastructure}.
\\ 
\hline

\textit{Attendance}
&Informasi tentang presensi setiap \textit{user} yang akan digunakan untuk perhitungan lembur.
&Data \textit{Attendance} akan ditambahkan ketika \textit{staff Infrastructure} melakukan presensi.
\\ 
\hline

\textit{Overtime}
&Informasi tentang lembur dari setiap \textit{staff Infrastructure}.
&Data \textit{Overtime} akan ditambahkan secara otomatis jika \textit{staff Infrastructure} melakukan presensi diluar jam \textit{shift} ataupun pada hari libur. Data \textit{Overtime} juga dapat ditambahkan secara manual oleh \textit{staff Infrastructure}.
\\ 
\hline

\textit{Session}
&Informasi tentang \textit{Session} setiap \textit{user}.
&Data \textit{Session} akan diregenerasi secara otomatis ketika \textit{user} melakukan \textit{login}.
\\
\hline

\textit{ActivityType}
&Informasi tentang jenis kerja yang dilakukan oleh \textit{staff} seperti WFO, WFH atau \textit{remote}.
&Data \textit{ActivityType} ditambahkan secara \textit{default} dengan nilai yang telah ditentukan.
\\
\hline

\textit{ActivityCategory}
&Informasi tentang kategori aktivitas pekerjaan yang dilakukan oleh \textit{staff Infrastructure} seperti instalasi dan sebagainya.
&Data \textit{ActivityCategory} ditambahkan secara \textit{default} dengan nilai yang telah ditentukan.
\\
\hline

\textit{ShiftScheduling}
&Informasi tentang penjadwalan \textit{shift} yang akan digunakan oleh setiap \textit{staff} berdasarkan jangka waktu yang diatur.
&Data \textit{ShiftScheduling} ditambahkan oleh \textit{Manager Human Resources} dan \textit{Manager Infrastructure}.
\\
\hline

\end{longtable}


%%%%%%%%%%%%%%%%%%%%%%%%%%%%%%%%%%%%%%%%%%%%%%%%%%%%%%%%%%%%%%%%%
\FloatBarrier
\begin{longtable}{ |p{2cm}|p{3cm}|p{3cm}|p{3cm}|p{3cm}|}
\caption{Identifikasi Hubungan Antar \textit{Entity}}\label{tab:indentifikasi-hubungan-antar-entity} 
\\
\hline
\textbf{\textit{Entity}}
&\textbf{\textit{Multiplicity}}
&\textbf{\textit{Relationship}}
&\textbf{\textit{Multiplicity}}
&\textbf{\textit{Entity}}
\\ 
\hline

\multirow{2}{*}{\textit{Shift}}
&1..1
&\textit{AssignedOn}
&1..*
&\textit{ShiftDay}
\\ 
\cline{2-5}
&1..1
&\textit{ShiftAssignments}
&0..*
&\textit{ShiftScheduling}
\\ 
\hline

\multirow{3}{*}{\textit{Attendance}}
&0..1
&\textit{Includes}
&0..1
&\textit{Overtime}
\\
\cline{2-5}
&1..*
&\textit{CategorizedBy}
&1..1
&\textit{ActivityType}
\\ 
\cline{2-5}
&1..*
&\textit{ClassifiedAs}
&1..1
&\textit{ActivityCategory}
\\ 
\hline

\multirow{8}{*}{\textit{User}}
&1..*
&\textit{WorksIn}
&1..1
&\textit{Department}
\\
\cline{2-5}
&1..*
&\textit{AssignedAs}
&1..1
&\textit{Role}
\\ 
\cline{2-5}
&1..1
&\textit{Records}
&0..*
&\textit{Attendance}
\\ 
\cline{2-5}
&0..1
&\textit{Requests}
&0..*
&\textit{Overtime}
\\ 
\cline{2-5}
&1..*
&\textit{ScheduledFor}
&0..1
&\textit{Shift}
\\ 
\cline{2-5}
&1..1
&\textit{Attends}
&1..1
&\textit{Session}
\\ 
\cline{2-5}
&1..*
&\textit{EntitledTo}
&0..1
&\textit{Holiday}
\\ 
\cline{2-5}
&1..1
&\textit{HasSchedules}
&0..*
&\textit{ShiftScheduling}
\\ 
\hline

\textit{Holiday}
&1..1
&\textit{AppliesOn}
&1..*
&\textit{HolidayDay}
\\ 
\hline


\end{longtable}

\subsection{\textit{Logical Database Design}}
\textit{Logical data modeling} yaitu proses menerjemahkan representasi konseptual ke dalam struktur logis basis data, yang mencakup perancangan relasi~\cite{conolly2015database}. Pada bagian \textit{logical data modeling} akan menjelaskan atribut dari setiap entitas. Diagram rancangan \textit{logical data modeling} dapat dilihat pada \hyperref[sec:logical-database-design]{\textbf{Lampiran~IX}}.

\subsection{\textit{Table Specification}}
Pada \textit{table specification} akan dijelaskan tentang tipe data dari setiap entitas, deskripsi daris setiap atribut, dan menjelaskan setiap atribut apakah memiliki status \textit{nullable, composite}, atau \textit{multivalue}. \textit{Table specification} dari \textit{database} aplikasi presensi \textit{online} pada PT Password Solusi Sistem dapat dilihat pada \hyperref[sec:table-specification]{\textbf{Lampiran~X}}.

